Sliding windows provide persistent-homology features, but shared observations can induce dependence and invalidate the exchangeability underlying window-level random permutation. This study evaluates the finite-sample calibration of a change-point test combining a fixed persistence-diagram representation (BPI) with a maximum scan statistic. A window-mean counterexample shows that i.i.d.\ raw observations need not yield exchangeable window features. Unrestricted window-level permutation (M0) and a dependent multiplier bootstrap (M1) are compared across 28 null configurations under fixed-window-count and fixed-observation-length designs, with mean--covariance features (MOMENTS) as a comparator. At the nominal level of 0.05, M0 is materially anti-conservative in all 20 overlapping configurations for each representation; the BPI rejection rate reaches 0.9756. M1 lowers rejection rates in all 56 matched configuration--representation combinations, yet is materially conservative in all eight non-overlapping configurations for each representation. For BPI with i.i.d.\ raw observations, equal averaging over the two generators yields M0 rejection-rate differences between strides 16 and 64 of 0.5354 and 0.5330 under the two designs. With 192 non-overlapping windows, the M1 rejection-rate deficit relative to 0.05 is 0.0266. Under a circle-to-figure-eight alternative with matched first and second marginal moments, size-adjusted powers are 1.0000 for BPI and 0.0620 for MOMENTS. A mean-shift negative control reverses the ordering, consistent with the translation invariance of persistent homology from Euclidean distances. Calibration is therefore condition-dependent for the studied implementation: lower rejection rates do not necessarily imply nominal calibration, and representation sensitivity and Type-I error control require separate evaluation.

\noindent\textbf{Keywords}: persistent homology; change-point detection; overlapping sliding windows; exchangeability; finite-sample calibration
